\documentclass[11pt,a4paper]{article}

\usepackage{amsmath,amssymb,amsthm}
\usepackage[margin=2.5cm]{geometry}
\usepackage{xcolor}
\usepackage{booktabs}
\usepackage{tabularx}
\usepackage{adjustbox}
\usepackage{enumitem}
\usepackage{microtype}

\usepackage[
  colorlinks=true,
  linkcolor=blue!55!black,
  citecolor=darkgray,
  urlcolor=blue!55!black,
  bookmarks=true,
  bookmarksnumbered=true,
  unicode=true
]{hyperref}

\geometry{a4paper, top=2.5cm, bottom=2.5cm, left=2.5cm, right=2.5cm}

\widowpenalty=10000
\clubpenalty=10000

\theoremstyle{plain}
\newtheorem{theorem}{Theorem}[section]
\newtheorem{lemma}[theorem]{Lemma}
\newtheorem{corollary}[theorem]{Corollary}
\newtheorem{proposition}[theorem]{Proposition}
\theoremstyle{definition}
\newtheorem{definition}[theorem]{Definition}
\newtheorem{example}[theorem]{Example}
\theoremstyle{remark}
\newtheorem{remark}[theorem]{Remark}

\newcommand{\norm}[1]{\lVert #1 \rVert}
\newcommand{\gap}{\operatorname{gap}}
\newcommand{\R}{\mathbb{R}}
\newcommand{\Z}{\mathbb{Z}}

\hypersetup{
  pdftitle={The Rung Family of Lonely Runner Spectra: gap(1,2,...,n-1,kn) = k/(kn+1)},
  pdfauthor={MIKEAA2020},
  pdfsubject={Lonely Runner Conjecture, lonely runner spectrum, exact gap computation},
  pdfkeywords={lonely runner, gap, rung family, argmax, filler characterization, census}
}

\begin{document}

\tableofcontents
\newpage

\section{Introduction}\label{sec:intro}

For a finite set $V$ of positive integer speeds, the \emph{loneliness} of $V$
is the quantity
\[
  \gap(V) \;=\; \sup_{t \in \R}\, \min_{v \in V} \norm{v t},
\]
where $\norm{x}$ denotes the distance of $x$ to the nearest integer. The
Lonely Runner Conjecture (LRC) asserts that every $m$-element set of distinct
speeds satisfies $\gap(V) \ge \tfrac{1}{m+1}$, and the bound is tight for the
consecutive family $V = \{1, 2, \dots, m\}$, for which
$\gap(V) = \tfrac{1}{m+1}$ exactly. The conjecture is known for up to seven
speeds and open from eight onward; exhaustive computational verification is
therefore one of the few systematic tools available, and the structure of the
\emph{spectrum} --- the multiset of gap values realized by all $m$-sets up to
some radius --- is a recurring theme of the computational literature.

This note is a self-contained study of one spectral family. For $n \ge 2$ and
$k \ge 1$ define the \emph{rung family}
\[
  F(n,k) \;=\; \{1, 2, \dots, n-1\} \cup \{kn\},
\]
an $n$-element set consisting of the first $n-1$ consecutive speeds together
with the single large speed $kn$. The rung $k=1$ is the classical consecutive
family. The values
\[
  \tfrac{1}{n+1},\ \tfrac{2}{2n+1},\ \tfrac{3}{3n+1},\ \dots,\ \tfrac{k}{kn+1}
\]
appear as a rigid arithmetic progression at the very bottom of every computed
spectrum --- the ``rungs'' of a ladder --- and the purpose of this note is to
prove that they are exact, to determine the complete maximizer structure
behind them, and to show how the same structure governs the neighbouring zoo
of extremal sets.

Three theorems are proved. \textbf{Theorem~\ref{thm:rung}} (Section
\ref{sec:value}) gives, for every $n \ge 2$ and $k \ge 1$, the exact value
$\gap(F(n,k)) = \tfrac{k}{kn+1}$ together with the \emph{complete} argmax
structure: for $k \ge 2$ the maximum is attained at the single pair
$t = \pm \tfrac{k}{kn+1} \pmod 1$, the binding runners are exactly
$\{1, kn\}$, and every other runner sits at distance at least
$\tfrac{k+1}{kn+1}$; for $k = 1$ the maximizers are exactly the unit
fractions $a/(n+1)$, a structural contrast that explains why the consecutive
family is exceptional. The upper bound is an elementary pigeonhole argument
that generalizes the proof of the $k=2$ case in the project's earlier audit;
the equality case is new in one decisive respect: the uniqueness of the
maximizer is proved by a \emph{no-wrap induction} that requires no
coprimality hypothesis, closing the gcd-degenerate case that the earlier
sum-based argument left open. \textbf{Theorem~\ref{thm:filler}} (Section
\ref{sec:filler}) is a closed-form characterization of the \emph{fillers} of
a core: for a speed set $S$ whose gap exceeds a rational target $g$, the set
$X(S) = \{x : \gap(S \cup \{x\}) = g\}$ is characterized exactly by two
finitely checkable conditions --- a window condition on each escape interval
of the core and a cleanliness condition at one level time --- and is finite
with an explicit a priori radius. This turns the empirically observed
ladder structure of extremal sets into a theorem. Finally, Section
\ref{sec:census} reports an exhaustive census of all
$\binom{50}{8} = 536{,}878{,}650$ eight-element speed sets with speeds up to
$50$, carried out with a filtered exact solver whose certificate is a
one-line consequence of the definition of $\gap$: the census shows the tight
zoo closed, the rung-2 zoo closed in primitive structure, and --- the target
question of the hunt --- the inter-rung window $\bigl(\tfrac19,
\tfrac{2}{17}\bigr)$ \emph{empty}.

All computations use exact integer or rational arithmetic throughout; no
floating point is involved anywhere. The machine checks summarized in
Section~\ref{sec:machine} were run with two independent implementations of
the gap computation that must agree on every value.

\section{Definitions and statements}\label{sec:statements}

Throughout, $\norm{x} = \min_{m \in \Z} |x - m|$ is the distance to the
nearest integer, so that $\norm{x} \in [0, \tfrac12]$ and
$\norm{-x} = \norm{x}$. Two elementary facts are used repeatedly:
$\norm{qx} \le q\,\norm{x}$ for every integer $q \ge 1$; and for a rational
$t = a/d$ in lowest terms, $\norm{v \cdot a/d}$ is computed exactly as
$\min(r, d-r)/d$ with $r = av \bmod d$.

\begin{definition}[gap, argmax, binders]\label{def:gap}
For a finite set $V$ of positive integers,
\[
  \gap(V) \;=\; \sup_{t \in \R}\, \min_{v \in V} \norm{vt},
\]
where the supremum may be taken over $t \in [0,1]$. The \emph{argmax set}
$\operatorname{Argmax}(V)$ is the set of $t \in [0,1)$ attaining the
supremum. For $t$ in the argmax set, the \emph{binders at $t$} are the
speeds $v \in V$ with $\norm{vt} = \gap(V)$.
\end{definition}

The function $g_V(t) = \min_{v \in V} \norm{vt}$ is continuous and piecewise
linear, with breakpoints at the peaks $t = \text{odd}/(2v)$ and crossings
$t(v_i \pm v_j) \in \Z$; its maximum is attained among these, which is what
the exact solvers enumerate. The Lonely Runner Conjecture in this
normalization states $\gap(V) \ge \tfrac{1}{|V|+1}$ for $|V|$ distinct
speeds, with equality for the consecutive family.

The three main results are as follows.

\begin{theorem}[rung family, value and complete argmax structure]
\label{thm:rung}
Let $n \ge 2$ and $k \ge 2$, and put $D = kn + 1$ and $\delta = k/D$.
Then:
\begin{enumerate}[label=\rm(\roman*), leftmargin=*]
\item $\gap\bigl(\{1, \dots, n-1, kn\}\bigr) = \dfrac{k}{kn+1}$;
\item $\operatorname{Argmax} = \Bigl\{ \dfrac{k}{kn+1},\ %
     \dfrac{kn+1-k}{kn+1} \Bigr\}$, i.e.\ the maximum is attained at
     exactly two times, reflections of each other;
\item at $t = \delta$ the binders are exactly the runners $1$ and $kn$,
     both at distance $\delta$;
\item every other runner $j$, $2 \le j \le n-1$, has
     $\norm{j\delta} \ge \dfrac{k+1}{kn+1} > \delta$.
\end{enumerate}
No coprimality condition on $(n,k)$ is required.
\end{theorem}

\begin{theorem}[base rung $k=1$]\label{thm:base}
For $n \ge 2$,
$\gap(\{1, 2, \dots, n\}) = \tfrac{1}{n+1}$, and the argmax set is exactly
\[
  \bigl\{\, a/(n+1) \;:\; 0 \le a \le n,\ \gcd(a, n+1) = 1 \,\bigr\},
\]
the grid of \emph{units} modulo $n+1$; at each such time $t = a/(n+1)$ the
binders are the speeds $j$ with $ja \equiv \pm 1 \pmod{n+1}$.
\end{theorem}

Thus the rung family exhibits a structural dichotomy: the base rung has
$\varphi(n+1)$ maximizers spread over the units grid, while every higher
rung collapses to a single maximizer pair. The proof of
Theorem~\ref{thm:rung} will show that the collapse is caused by exactly one
inequality, $D \ge 2n+1$, which prevents the residue walk of the equality
configuration from wrapping around the modulus.

\begin{theorem}[closed-form filler characterization]\label{thm:filler}
Let $S$ be a finite set of positive integers, $g \in (0,\tfrac12)$ rational,
and $g_S(t) = \min_{s \in S} \norm{st}$. Define
\[
  E(S) = \{ t \in [0,1] : g_S(t) > g \}, \qquad
  L(S) = \{ t \in [0,1] : g_S(t) = g \},
\]
the \emph{escape set} and the \emph{level set} of the core. Then $E(S)$ is a
finite union of open intervals $(\alpha_j, \beta_j)$ with $0 < \alpha_j <
\beta_j < 1$ and $g_S(\alpha_j) = g_S(\beta_j) = g$, and $L(S)$ is a finite
set containing $\{\alpha_j, \beta_j\}_j$. For an integer $x \ge 1$ with
$x \notin S$,
\[
  \gap(S \cup \{x\}) = g
  \iff
  \text{\rm(A)} \wedge \text{\rm(B)},
\]
where
\begin{itemize}[leftmargin=*]
\item[\text{\rm(A)}] for every escape interval $(\alpha_j, \beta_j)$ there is
  an integer $m_j$ with $m_j - g \le x\,\alpha_j$ and $x\,\beta_j \le m_j + g$
  (equivalently, $x(\beta_j - \alpha_j) \le 2g$ and the interval
  $[\,x\beta_j - g,\ x\alpha_j + g\,]$ contains an integer);
\item[\text{\rm(B)}] some $t \in L(S)$ satisfies $\norm{xt} \ge g$.
\end{itemize}
Moreover, if $\gap(S) > g$ then \text{\rm(A)} forces
$x \le 2g / \max_j (\beta_j - \alpha_j)$, so the filler set
\[
  X(S) := \{ x \ge 1 : \gap(S \cup \{x\}) = g \}
\]
is finite with the explicit radius $R(S) = \bigl\lfloor 2g /
\max_j(\beta_j - \alpha_j) \bigr\rfloor$; if $\gap(S) < g$ then $X(S) =
\emptyset$; and if $\gap(S) = g$ then $X(S) = \{x : \text{\rm(B)}\}$ is
infinite.
\end{theorem}

Condition (A) says that the filler must \emph{kill} every escape of the core:
wherever the core exceeds the target, the filler's sawtooth must be at most
$g$, which forces the image $x \cdot (\alpha_j, \beta_j)$ to lie inside a
single window $[m_j - g, m_j + g]$ of length $2g$ on the circle $\R/\Z$.
Condition (B) says the filler must be \emph{clean} (at distance $\ge g$) at
some level time of the core, which is precisely where the extension attains
the value $g$: a filler can bind only at a level time of the core. Together
the two conditions reduce the search for all fillers to a finite, exactly
checkable list of two-sided congruence inequalities, computed once from the
core's own geometry.

The proofs occupy Sections~\ref{sec:value}--\ref{sec:filler}, the census of
Section~\ref{sec:census} supplies the computational picture at $n = 8$, and
Section~\ref{sec:machine} summarizes the machine verification of every
statement in exact arithmetic.

\section{The upper bound}\label{sec:value}

\begin{lemma}[pigeonhole step]\label{lem:pigeon}
Let $n \ge 2$, $k \ge 2$, $D = kn+1$, $\delta = k/D$. For every real $t$,
\[
  \min\bigl( \norm{t}, \norm{2t}, \dots, \norm{(n-1)t}, \norm{knt} \bigr)
  \;\le\; \delta .
\]
\end{lemma}

\begin{proof}
Suppose, for contradiction, that $\norm{jt} > \delta$ for $1 \le j \le n-1$
and $\norm{knt} > \delta$, and consider the $n+1$ points
$P = \{0, t, 2t, \dots, nt\}$ on the circle $\R/\Z$.

\emph{Classical step.} Among $n+1$ points on the unit circle, two lie at
circular distance at most $\tfrac1{n+1}$. Their index difference $d$ lies in
$\{1, \dots, n\}$, so $\norm{dt} \le \tfrac1{n+1}$. Since
\[
  \tfrac{1}{n+1} < \tfrac{k}{kn+1} = \delta
  \iff k(n+1) > kn + 1 \iff k > 1,
\]
and $\norm{jt} > \delta$ for all $j \le n-1$, the difference must be $d = n$:
writing $\beta = \norm{nt}$,
\[
  \beta \;\le\; \tfrac{1}{n+1} \;<\; \delta .
\]

\emph{Adjacency.} For $0 \le i < j \le n$ with $(i,j) \ne (0,n)$, the index
difference $j - i$ lies in $\{1, \dots, n\}$ and, whenever $j - i \le n-1$,
is the index of a runner, so the circular distance of $it$ and $jt$ equals
$\norm{(j-i)t} > \delta$. The only pair with $j - i = n$ is $(0,n)$ itself.
Hence $\{0, nt\}$ is the unique pair at distance below $\delta$, and $0$ and
$nt$ are adjacent in the circular order: a point $z$ strictly inside the
short arc from $0$ to $nt$ would satisfy $\norm{z} < \beta < \delta$, but
$z = jt$ with $1 \le j \le n-1$, a contradiction.

\emph{Gap counting.} The $n+1$ circular gaps of $P$ consist of the gap
$\beta$ (between $0$ and $nt$) and $n$ further gaps, each bounded below by
the circular distance of its endpoint pair, which is a pair other than
$\{0, nt\}$ and hence has distance $> \delta$. Summing to $1$,
\[
  1 \;>\; \beta + n\delta \;=\; \beta + \frac{nk}{kn+1}
  \quad\Longrightarrow\quad
  \beta \;<\; \frac{1}{kn+1} \;=\; \frac{\delta}{k}.
\]

\emph{Kill.} Finally $\norm{knt} = \norm{k \cdot nt} \le k\,\norm{nt} =
k\beta < k \cdot \tfrac{\delta}{k} = \delta$, contradicting
$\norm{knt} > \delta$.
\end{proof}

The lower bound half of Theorem~\ref{thm:rung}(i) is a direct evaluation at
$t = \delta$: $\norm{1 \cdot \delta} = \delta$, and
\[
  kn \cdot k \;=\; k(kn) \;\equiv\; k \cdot (-1) \;=\; -k \pmod{D},
\]
since $kn \equiv -1 \pmod{D}$, whence $\norm{kn\,\delta} = k/D = \delta$
(because $2k < kn+1$ for $n \ge 2$). So $\gap(F(n,k)) \ge \delta$, and
Lemma~\ref{lem:pigeon} gives equality. This proves
Theorem~\ref{thm:rung}(i) for all $k \ge 2$. For $k = 1$ the classical step
already suffices: $d \in \{1, \dots, n\}$ are \emph{all} runners when $k=1$
(the speed $n$ is present), so $\norm{dt} \le \tfrac1{n+1}$ directly; and
$t = \tfrac1{n+1}$ attains $\tfrac1{n+1}$ through runners $1$ and $n$. This
is the value half of Theorem~\ref{thm:base}.

\section{The equality case: no-wrap induction}\label{sec:equality}

Throughout this section $n \ge 2$, $k \ge 2$, $D = kn+1$, $\delta = k/D$,
and $t$ is a time at which the minimum equals $\delta$, i.e.\
$\norm{jt} \ge \delta$ for $1 \le j \le n-1$ and $\norm{knt} \ge \delta$.
The analysis has four steps, and it is the fourth that removes the
coprimality hypothesis used in earlier treatments.

\begin{lemma}[gap multiset]\label{lem:multiset}
In the equality case, the circular gaps of $P = \{0, t, \dots, nt\}$ are
exactly $\{\tfrac1D,\ \tfrac{k}{D}, \dots, \tfrac{k}{D}\}$ ($n$ copies), the
gap of length $\tfrac1D$ lying between $0$ and $nt$.
\end{lemma}

\begin{proof}
Repeating the classical step with the weak hypothesis gives
$\beta := \norm{nt} \le \tfrac1{n+1} < \delta$, so $\{0, nt\}$ is still the
unique pair below $\delta$ and $0, nt$ are adjacent. The other $n$ gaps are
each at least the circular distance of their endpoint pair, which is a
runner difference of index in $\{1, \dots, n-1\}$, hence $\ge \delta$. If
$\beta < \tfrac1D$, then $\norm{knt} \le k\beta < \tfrac{k}{D} = \delta$
contradicts the hypothesis; so $\beta = \tfrac1D$. The $n$ remaining gaps
then sum to $1 - \tfrac1D = \tfrac{nk}{D} = n\delta$, each being $\ge \delta$,
so all equal $\delta = \tfrac{k}{D}$.
\end{proof}

Scaling by $D$, the points of $P$ in circular order are partial sums of a
sequence of $n+1$ steps consisting of $n$ steps of size $k$ and one step of
size $1$ (the $\beta$-step between $nt$ and $0$). If the unit step occupies
position $p+1$ in the circular order, $0 \le p \le n$, then
\begin{equation}\label{eq:Ap}
  P \;=\; \Bigl\{ \frac{jk}{D} : 0 \le j \le p \Bigr\}
  \;\cup\;
  \Bigl\{ \frac{jk+1}{D} : p \le j \le n-1 \Bigr\},
\end{equation}
the case $p = n$ being the pure arithmetic progression and $p = 0$ the
reflected one.

\begin{lemma}[no mixed configurations]\label{lem:mixed}
$p \in \{0, n\}$.
\end{lemma}

\begin{proof}
Suppose $1 \le p \le n-1$. Then both $\tfrac{pk}{D}$ and
$\tfrac{pk+1}{D}$ belong to $P$ by \eqref{eq:Ap}, and they are distinct
points at circular distance $\tfrac1D < \delta$. Being two points of $P$,
they are $it$ and $jt$ for some indices $0 \le i, j \le n$, whose index
difference lies in $\{1, \dots, n\}$. If $|i - j| \le n - 1$, then
$\norm{(i-j)t} = \tfrac1D < \delta$ contradicts $\norm{(i-j)t} \ge \delta$
(the difference is a runner). Hence $|i-j| = n$, so $\{i, j\} = \{0, n\}$
and the two points are $0$ and $nt$. But $0 < pk < D$ and
$0 < pk + 1 \le (n-1)k + 1 < D$, so neither $pk$ nor $pk+1$ is $0$ modulo
$D$: neither point is $0$, a contradiction.
\end{proof}

\begin{lemma}[no-wrap induction]\label{lem:nowrap}
Let $D = kn+1$ with $k \ge 2$, and let $m \in \{1, \dots, n\}$ satisfy
\begin{equation}\label{eq:seteq}
  \{ jm \bmod D : 0 \le j \le n \} \;=\; \{0, 1, \dots, n\}.
\end{equation}
Then $m = 1$.
\end{lemma}

\begin{proof}
Write $\sigma(j) = jm \bmod D \in \{0, \dots, n\}$ for $0 \le j \le n$;
the hypothesis \eqref{eq:seteq} says precisely that these are the residues of
the walk and that all of them land in $\{0, \dots, n\}$. Now $\sigma(0) = 0$,
and for each $j$,
\[
  \sigma(j+1) \;\equiv\; \sigma(j) + m \pmod D .
\]
Both $\sigma(j+1)$ and $\sigma(j) + m$ lie in $\{0, \dots, D-1\}$: the
former by definition, the latter because
\[
  \sigma(j) + m \;\le\; n + n \;=\; 2n \;\le\; D - 1,
\]
using $D = kn + 1 \ge 2n + 1$. Two numbers in $\{0, \dots, D-1\}$ that are
congruent modulo $D$ are equal, so $\sigma(j+1) = \sigma(j) + m$
\emph{without reduction}, for every $j$. By induction $\sigma(j) = jm$ as
integers, so $jm \le n$ for all $j \le n$; taking $j = n$ gives $nm \le n$,
i.e.\ $m \le 1$, and $m \ge 1$ finishes the proof.
\end{proof}

\begin{proof}[Proof of Theorem~\ref{thm:rung}]
Part (i) was proved at the end of Section~\ref{sec:value}. For (ii)--(iv),
let $t$ be a time attaining the minimum $\delta$. By
Lemmas~\ref{lem:multiset} and \ref{lem:mixed}, $P$ is either the pure
progression $\{\tfrac{jk}{D} : 0 \le j \le n\}$ ($p = n$) or its reflection
($p = 0$). In the pure case $t \equiv \tfrac{mk}{D} \pmod 1$ for some
$m \in \{1, \dots, n\}$ (as $t$ is itself a point of $P$), and the point
set equality $\{jt \bmod 1\} = P$ reads
\[
  \{ jmk \bmod D : 0 \le j \le n \} \;=\; \{ jk \bmod D : 0 \le j \le n \};
\]
multiplying by the unit $k^{-1} \bmod D$ (note $\gcd(k, kn+1) = 1$) yields
\eqref{eq:seteq} of Lemma~\ref{lem:nowrap}, so $m = 1$ and $t =
\tfrac{k}{D} = \delta$. In the reflected case, $P = \{\tfrac{jk+1}{D} : 0
\le j \le n-1\}$; applying the pure case to $1 - t$ (which has the same gap
value, since $\norm{j(1-t)} = \norm{jt}$, and whose point set is
$\{j(1-t)\} = -\{jt\} = \{\tfrac{jk}{D} : 1 \le j \le n\} \cup \{0\}$, the
pure progression) yields $1 - t = \delta$, i.e.\ $t = 1 - \delta =
\tfrac{D-k}{D}$. This proves (ii).

For (iii) and (iv), evaluate at $t = \delta$: $\norm{\delta} = \delta$
(runner $1$); $\norm{kn\,\delta} = \norm{-k/D} = \delta$ (runner $kn$);
and for $2 \le j \le n-1$,
\[
  \norm{j\delta} \;=\; \frac{\min\bigl(jk,\; D - jk\bigr)}{D}
  \;\ge\; \frac{\min\bigl(2k,\; k+1\bigr)}{D}
  \;=\; \frac{k+1}{kn+1} \;>\; \delta,
\]
since $jk \ge 2k \ge k+1$ and $D - jk \ge D - (n-1)k = k+1$.
\end{proof}

\begin{remark}[why the gcd condition disappears]\label{rem:gcd}
The earlier treatment of the $k = 2$ case --- and its immediate
generalization --- derived \eqref{eq:seteq} and then summed the least
residues, obtaining $(m-1)\,n(n+1)/2 = (kn+1)\,Q$; concluding $m = 1$
required $\gcd\bigl(n(n+1)/2,\ kn+1\bigr) = 1$, which always holds for
$k = 2$ but can fail for $k \ge 3$ (e.g.\ $n = k = 3$, where the gcd is
$2$). Lemma~\ref{lem:nowrap} replaces the sum argument by the observation
that the residue walk $\sigma$ never wraps, which is where $k \ge 2$ --- and
only that --- is used. As a complement, the arithmetic obstruction itself is
harmless: since $\gcd(n, kn+1) = 1$, one always has
$\gcd\bigl(n(n+1)/2,\ kn+1\bigr)$ dividing $n+1$, hence dividing
$k(n+1) - (kn+1) = k - 1$; so the degenerate modulus can only be a divisor
of $k-1$, and for $k \in \{2,3\}$ the sum argument already suffices for all
$n$. The no-wrap induction covers all $k \ge 2$ uniformly.
\end{remark}

\begin{remark}[the base rung is different]\label{rem:base-contrast}
For $k = 1$ one has $D = n + 1 \le 2n$, so the inequality
$\sigma(j) + m \le D - 1$ fails --- and indeed the conclusion fails:
by Theorem~\ref{thm:base} there are $\varphi(n+1)$ maximizers, not two.
The rung family thus exhibits a sharp dichotomy at $k = 1$ versus $k \ge 2$,
governed precisely by whether $D \ge 2n+1$.
\end{remark}

\section{The base rung}\label{sec:base}

\begin{proof}[Proof of Theorem~\ref{thm:base}]
The value was proved in Section~\ref{sec:value}. Let $t$ attain the minimum
$\tfrac{1}{n+1} =: \delta$ for the speed set $\{1, \dots, n\}$, so
$\norm{jt} \ge \delta$ for all $1 \le j \le n$. The classical step gives
$\beta := \norm{nt} \le \tfrac1{n+1} = \delta$, and $0, nt$ are adjacent. But
now runner $n$ is present in the speed set, so $\beta = \norm{nt} \ge
\delta$ at the argmax; hence $\beta = \delta$ exactly, and the $n$ other
gaps, each $\ge \delta$ and summing with $\beta$ to $1 = (n+1)\delta$, all
equal $\delta$ as well. So all $n+1$ circular gaps of $P$ equal $\delta =
\tfrac1{n+1}$: $P$ is the \emph{equally spaced} set $\{0, \tfrac1{n+1},
\tfrac2{n+1}, \dots, \tfrac{n}{n+1}\}$.

Hence $t = \tfrac{m}{n+1}$ for some $m \in \{1, \dots, n\}$, and the point
set equality $\{jt \bmod 1\} = P$ becomes
\[
  \{ jm \bmod (n+1) : 0 \le j \le n \} \;=\; \{0, 1, \dots, n\},
\]
which holds if and only if $m$ is a unit modulo $n+1$: if
$c = \gcd(m, n+1) > 1$, every residue $jm$ is divisible by $c$ and the
right-hand side cannot be covered. Conversely, for a unit $m$ the map $j
\mapsto jm$ permutes the residues. At $t = \tfrac{a}{n+1}$ with $a$ a unit,
the value is $\min_j \norm{\tfrac{ja}{n+1}} = \tfrac{1}{n+1}$, attained
exactly by the speeds $j$ with $ja \equiv \pm 1 \pmod{n+1}$ (which exist
since $a$ is invertible).
\end{proof}

The two theorems together give the rung ladder its exact shape. In every
computed spectrum the bottom values are precisely the rung values
$\tfrac{k}{kn+1}$ present at the given radius: for $n = 7$ the ladder reads
$\tfrac18, \tfrac2{15}, \tfrac3{22}, \tfrac4{29}, \tfrac5{36}, \tfrac6{43},
\tfrac7{50}$ and for $n = 8$ it reads $\tfrac19, \tfrac2{17}, \tfrac3{25},
\tfrac4{33}, \tfrac5{41}, \tfrac6{49}$ --- each realized by the family member
$F(n,k)$ itself, and, as Section~\ref{sec:census} shows, each carrying its
own zoo of further extremal sets.

\section{The filler characterization}\label{sec:filler}

The rung values sit at the bottom of the spectrum, and the sets attaining
them form the \emph{zoo}: at $n = 7$, radius $V = 50$, exactly $29$ sets have
gap $\tfrac{2}{15}$, and at $n = 8$, $V = 50$, exactly $21$ sets have gap
$\tfrac{2}{17}$. Every zoo member ever observed decomposes as a \emph{core}
(a set whose own gap exceeds the rung value) plus one \emph{filler} speed,
and the empirical classification of zoos ran through the question: which
fillers does a given core admit? Theorem~\ref{thm:filler} answers it in
closed form. This section proves the theorem and works the central example.

\subsection{Geometry of the core}

Fix the core $S$ and the rational target $g \in (0, \tfrac12)$, and write
$g_S(t) = \min_{s \in S} \norm{st}$, a continuous piecewise-linear function.
Each sawtooth $\norm{st}$ has breakpoints at $t = \text{odd}/(2s)$ and
crosses the level $g = p/q$ exactly at the rationals $t = (qm \pm p)/(sq)$.
Let $\tau_0 < \tau_1 < \dots < \tau_K$ be the sorted union of $\{0, 1\}$,
all breakpoints and all level crossings of all sawtooths. On each open
segment $(\tau_i, \tau_{i+1})$ every $\norm{st}$ is linear and does not
cross level $g$, so the sign of $g_S - g$ is constant there. The escape set
$E = \{g_S > g\}$ is therefore a finite union of maximal open intervals
$(\alpha_j, \beta_j)$; merging across an event point $\tau$ is legitimate
only when $g_S(\tau) > g$, so each $\alpha_j, \beta_j$ carries
$g_S = g$ exactly, i.e.\ $\alpha_j, \beta_j \in L(S)$. Neither endpoint can
be $0$ or $1$, because $g_S(t) \to 0$ as $t \to 0$ or $1$; the level set
$L(S)$ is finite, being contained in the crossing set of the individual
sawtooths (no $\norm{st}$ is constant on an interval unless two distinct
speeds $s \ne s'$ have $\norm{st} = \norm{s't} = g$ on an interval, which
forces $(s - s')t \in \Z$ on an interval and hence $s = s'$).

\begin{proof}[Proof of Theorem~\ref{thm:filler}]
Put $h(t) = \min\bigl(g_S(t), \norm{xt}\bigr)$, so that $\gap(S \cup \{x\})
= \max_{[0,1]} h$.

\emph{The bound $\gap \le g$ is condition (A).} Outside $E$ the core already
has $g_S \le g$. On an escape interval $(\alpha, \beta)$ the filler must
carry the bound: $\norm{xt} \le g$ for all $t \in (\alpha, \beta)$. Writing
the circle $\R/\Z$ as the union of the closed windows $[m - g, m + g]$,
$m \in \Z$, of length $2g$ separated by open gaps of length $1 - 2g > 0$,
the image interval $[x\alpha, x\beta]$ (for $x > 0$) is covered by the
windows if and only if it lies inside a single window $[m_j - g, m_j + g]$:
a connected interval cannot straddle a gap. By continuity the same holds on
the closed $[\alpha, \beta]$, and this is exactly (A).

\emph{Equality needs a touch (B).} Assume (A), so $\gap \le g$. The maximum
of the continuous $h$ is attained at some $t_0$. If $h(t_0) = g$, then
$g_S(t_0) \ge g$ and $\norm{xt_0} \ge g$. If $g_S(t_0) > g$ then $t_0 \in
E$; inside an escape interval, the window condition gives $\norm{xt_0} =
\operatorname{dist}(x t_0, m_j)$, and an interior point at window distance
exactly $g$ would force an endpoint past the window edge, contradicting
(A) --- formally, $x t_0 = m_j \pm g$ with $x\alpha \le x t_0 \le x\beta$
and $x\alpha, x\beta \in [m_j - g, m_j + g]$ forces $x t_0 \in
\{x\alpha, x\beta\}$, an endpoint --- and $\norm{xt} < g$ throughout the
open interval would give $h < g$ there. Hence an interior touch cannot
occur, the touch must be at an endpoint (where $g_S = g$), or
$g_S(t_0) = g$ holds outright. Either way $t_0 \in L(S)$ and
$\norm{xt_0} \ge g$: this is (B). Conversely, if $t \in L(S)$ has
$\norm{xt} \ge g$, then $h(t) = \min(g, \norm{xt}) = g$, so $\gap \ge g$,
and with (A) equality follows.

\emph{Degenerate cores.} If $\gap(S) < g$ then $L(S) = \emptyset$ and $h \le
g_S \le \gap(S) < g$: no fillers. If $\gap(S) = g$ then $E = \emptyset$
and (A) is vacuous; $\gap(S \cup \{x\}) = \max \min(g_S, \norm{x \cdot})$
equals $g$ exactly when some $t$ has $g_S(t) = g$ and $\norm{xt} \ge g$,
i.e.\ (B) alone; the solutions form infinite arithmetic families since
$\norm{xt} \ge g$ at a fixed rational $t$ excludes only finitely many
residue classes.

\emph{Radius.} If $E \ne \emptyset$, condition (A) applied to the longest
escape interval gives $x (\beta - \alpha) \le 2g$ for that interval, hence
$x \le 2g / \max_j (\beta_j - \alpha_j)$.
\end{proof}

\begin{example}[the hub $H_2$]\label{ex:H2}
Take $n = 7$, $g = \tfrac{2}{15}$, and the core
\[
  H_2 \;=\; (1, 4, 5, 6, 7, 11),
  \qquad \gap(H_2) = \tfrac{2}{13}.
\]
The escape set consists of two mirror intervals,
\[
  E(H_2) \;=\; \bigl(\tfrac{32}{105},\ \tfrac{14}{45}\bigr)
  \;\cup\;
  \bigl(\tfrac{31}{45},\ \tfrac{73}{105}\bigr),
  \qquad \beta - \alpha = \tfrac{2}{315},
\]
and the level set is
\[
  L(H_2) \;=\; \Bigl\{ \tfrac{32}{105},\ \tfrac{14}{45},\ \tfrac{7}{15},\
  \tfrac{8}{15},\ \tfrac{31}{45},\ \tfrac{73}{105} \Bigr\},
\]
of which $\tfrac{7}{15}$ and $\tfrac{8}{15}$ are \emph{isolated} level times
(local maxima of $g_{H_2}$ at height exactly $\tfrac{2}{15}$, where the
speeds $4$ and $11$ bind). The radius is $R = \bigl\lfloor
\tfrac{4/15}{2/315} \bigr\rfloor = 42$. Evaluating (A) on the escape
intervals leaves the candidates $\{3, 10, 13, 16, 26, 39\}$ from
$\{1, \dots, 42\}$; the touch (B) at $t = \tfrac{7}{15}$ requires
$\norm{7x/15} \ge \tfrac{2}{15}$, i.e.\ $7x \not\equiv 0, \pm 1 \pmod{15}$,
which prunes exactly $x = 13$ ($91 = 6 \cdot 15 + 1$). The theorem yields
\[
  X(H_2) \;=\; \{3,\ 10,\ 16,\ 26,\ 39\},
\]
matching brute force on $[1, 70]$ exactly, with no filler beyond the
provable radius $R = 42$: the five fillers generate the five zoo classes
$(1,3,4,5,6,7,11)$, $(1,4,5,6,7,10,11)$, $(1,4,5,6,7,11,16)$,
$(1,4,5,6,7,11,26)$, $(1,4,5,6,7,11,39)$ of the $n = 7$ census. Note that
$26 \equiv 11 \pmod{15}$ lands in the binder class itself: the extension by
$26$ binds with three speeds $\{4, 11, 26\}$ at $t^\ast = \tfrac{7}{15}$.
\end{example}

\begin{corollary}[fillers bind at level times]\label{cor:bindlevel}
If $\gap(S \cup \{x\}) = g$ and $t^\ast$ is an argmax, then $g_S(t^\ast) =
g$, i.e.\ $t^\ast \in L(S)$; in particular a filler never binds at a time
where the core is strictly above $g$ (the window argument forbids $\norm{x
 t} = g$ inside an escape interval), and whenever the filler binds, the core
is at level $g$ there as well.
\end{corollary}

This is the precise sense in which the theorem is ``Lemma A in reverse'':
the generalized binding-pair lemma states that an argmax at value $g =
\tfrac{2}{2n+1}$ lives on a grid $d = (2n+1)e$ with binders $v, w$ satisfying
$va \equiv \pm 2e \pmod d$ (so $v + w \equiv 0 \bmod d$); the level times of
a core at height $g$ are exactly where such residues can occur, and
condition (B) is the residue-cleanliness of the filler there.

\subsection{Verification}

Theorem~\ref{thm:filler} was verified by exact comparison against brute
force: for every core of every zoo member at $n = 5, 6, 7, 8$ (all $612$
distinct cores of size $|V| - 1$) and for $480$ random cores above the
rung, the theorem's $X(S)$ agrees with the brute-force filler list on
$[1, 70]$ with zero mismatches, using two independent implementations of
the gap computation (the sweep solver's integer routine and a Python
rational reference) plus a third cross-check. The degenerate regimes (cores
on and below the rung) behave as stated; below-rung cores never occur in
the zoo regimes, an immediate consequence of the LRC holding for up to
seven speeds, since a core of size $m$ has $\gap \ge \tfrac1{m+1} >
\tfrac{2}{2m+3}$.

\section{The \texorpdfstring{$n = 8$}{n=8} census at \texorpdfstring{$V = 50$}{V=50}}\label{sec:census}

The inter-rung question is the target of the hunt: \emph{does any $8$-element
speed set have gap strictly between $\tfrac19$ and $\tfrac{2}{17}$?} At $n =
6$ and $n = 7$ the corresponding windows are populated ($2$ and $4$ sets
respectively, the $\tfrac{3}{23}$ counterexamples), so $n = 8$ is the first
case where the window could close, and it was verified empty through $V =
38$ ($\binom{38}{8} = 48{,}903{,}492$ sets) in the project's earlier sweep.
This section settles the same radius as the $n = 7$ census.

\subsection{Method: the probe-lemma filter}

The census enumerates all $\binom{50}{8} = 536{,}878{,}650$ sets. The
feasible cost rests on a one-line certificate: for any fixed probe time
$t_0$,
\[
  \gap(V) \;\ge\; \min_{v \in V} \norm{v t_0},
\]
because the supremum over $t$ dominates the value at $t_0$. The hunt solver
(v4) scans the solver's own exact candidate times --- peaks
$\text{odd}/(2v)$ and crossings $t(v_i \pm v_j) \in \Z$, the complete
breakpoint set of $g_V$ --- and skips the vector as soon as one candidate
certifies a value above the threshold $\theta = \tfrac{2}{15}$, which
exceeds every hunted band (tight $\tfrac19$, inter-rung $\bigl(\tfrac19,
\tfrac{2}{17}\bigr)$, rung values $\tfrac{k}{8k+1}$ for $k \le 6$, all
$\le \theta$). Vectors with gap $\le \theta$ --- precisely the hunted
population --- always fall through to the full exact solve. Sampled vectors
(every $99{,}733$rd) are always fully solved, preserving sample validation.
The filter skipped $99.9985\%$ of vectors; the entire census ran in
$215$ seconds of single-core time.

Validation of the filtered census: (i) on the full $V = 38$ range, the
tight, rung-2 and inter-rung files and the exact bottom spectrum (all
$1{,}565$ vectors with gap $\le \tfrac{2}{15}$) are identical to the
unfiltered v3 solver's, and the shard-level outputs are byte-identical;
(ii) on a $V = 50$ sub-range, filtered and unfiltered runs agree
byte-for-byte on all census files and samples (a $105\times$ speedup on
that range); (iii) all $5{,}384$ sample records were re-verified against an
independent rational implementation; (iv) the accounting is exact:
$\text{spectrum} + \text{filtered} = 536{,}878{,}650$.

\subsection{Results}

\begin{enumerate}[label=\textbf{R\arabic*.}, leftmargin=*]
\item \textbf{Tight zoo closed.} $\gap = \tfrac19$ occurs exactly for
$\{1, \dots, 8\} \times \{1, 2, 3, 4, 5, 6\}$ --- the consecutive family and
its scalars, nothing else, at any radius up to $50$.
\item \textbf{Rung-2 zoo closed in primitive structure.} $\gap =
\tfrac{2}{17}$ occurs for $21$ vectors in exactly the same $7$ primitive
classes as at $V = 38$, with scalar multiples now up to $4$ (Table
\ref{tab:rung2}). The seven vectors with $v_8 > 38$ are precisely
$\times 2, \times 3, \times 4$ scalings of the existing classes: no new
binding structure appears.
\item \textbf{Inter-rung window empty.} No $8$-set with speeds up to $50$
has gap in $\bigl(\tfrac19, \tfrac{2}{17}\bigr)$. The spectrum jumps
directly from $\tfrac19$ ($6$ vectors) to $\tfrac{2}{17}$ ($21$ vectors);
the least attainable value strictly above $\tfrac19$ is the rung-2 value
itself.
\item \textbf{Every rung value carries a zoo.} The values
$\tfrac{3}{25}, \tfrac{4}{33}, \tfrac{5}{41}, \tfrac{6}{49}$ are attained by
$11$, $10$, $5$ and $3$ vectors respectively; the family members
$\{1, \dots, 7, 8k\}$ are present in each, and every member satisfies the
generalized binding-pair lemma on grids $d = (8k+1)e$ (Table~\ref{tab:rungk}).
The rung-5 and rung-6 zoos \emph{enter} at this radius, since they need
speeds $40$ and $48$.
\item \textbf{The ladder is certified by the filler theorem.} Each of the
seven new rung-2 members decomposes as core $+$ filler with the filler in
$X(\text{core})$ computed by Theorem~\ref{thm:filler} at $g =
\tfrac{2}{17}$; all eight $7$-subsets of each member support the
decomposition. For instance $(3, 6, 9, 12, 15, 18, 21, 48) = 3 \times
(1, 2, 3, 4, 5, 6, 7, 16)$ has core $(3, 6, 9, 12, 15, 18, 21)$ with
filler $48 \in X(\text{core})$, $R = 84$.
\item \textbf{Feasibility.} Every one of the $536{,}878{,}650$ vectors is
feasible ($\gap \ge \tfrac19$): the Lonely Runner Conjecture holds for
$8$ speeds up to $V = 50$, extending the $V = 38$ verification.
\end{enumerate}

\begin{table}[htbp]
\centering
\caption{The rung-2 zoo at $n = 8$, $V = 50$: $21$ vectors, $7$ primitive
classes (binding pairs modulo $17$; grids $d = 17e$ with $e \le 3$).}
\label{tab:rung2}
\small
\begin{tabular}{@{}lcc@{}}
\toprule
primitive class & scalars & binding pair \\
\midrule
$A = (1,2,3,4,5,6,7,16)$   & $\times\{1,2,3\}$ & $\{1,16\}$ \\
$B = (1,2,3,4,5,7,8,12)$   & $\times\{1,2,3,4\}$ & $\{5,12\}$ \\
$C = (1,2,3,4,5,7,12,16)$  & $\times\{1,2,3\}$ & $\{1,16\}$ \\
$D = (1,2,5,6,7,8,11,13)$  & $\times\{1,2,3\}$ & $\{6,11\}$ \\
$E = (1,4,5,6,7,8,11,13)$  & $\times\{1,2,3\}$ & $\{6,11\}$ \\
$F = (1,4,5,6,7,11,13,16)$ & $\times\{1,2,3\}$ & $\{1,16\}, \{4,13\}, \{6,11\}$ \\
$G = (1,5,6,7,8,11,13,21)$ & $\times\{1,2\}$ & $\{6,11\}, \{13,21\}$ \\
\bottomrule
\end{tabular}
\end{table}

\begin{table}[htbp]
\centering
\caption{Rung-$k$ zoos at $n = 8$, $V = 50$. Each zoo contains the family
member $\{1, \dots, 7, 8k\}$ and satisfies the binding-pair lemma on grids
$d = (8k+1)e$.}
\label{tab:rungk}
\small
\begin{tabular}{@{}ccccc@{}}
\toprule
$k$ & value & vectors & primitive classes & family member present \\
\midrule
1 & $1/9$    & 6  & 1 & yes (consecutive $\times\,1..6$) \\
2 & $2/17$   & 21 & 7 & yes \\
3 & $3/25$   & 11 & 6 & yes \\
4 & $4/33$   & 10 & 10 & yes \\
5 & $5/41$   & 5  & 5 & yes (new at this radius) \\
6 & $6/49$   & 3  & 3 & yes (new at this radius) \\
\midrule
\multicolumn{2}{l}{inter-rung $(1/9, 2/17)$} & \multicolumn{3}{l}{\textbf{0} (empty at $V \le 50$)} \\
\bottomrule
\end{tabular}
\end{table}

The bottom of the exact spectrum (all values $\le \tfrac{2}{15}$, complete
at $V = 50$) reads
\begin{center}\small
$\begin{array}[t]{lccccccc}
\text{value} & \tfrac19 & \tfrac2{17} & \tfrac3{25} & \tfrac4{33} &
\tfrac5{41} & \tfrac6{49} & \tfrac18 \\[2pt]
\text{count} & 6 & 21 & 11 & 10 & 5 & 3 & 763
\end{array}$\qquad
$\begin{array}[t]{lccccccc}
\text{value} & \tfrac7{55} & \tfrac6{47} & \tfrac5{39} & \tfrac4{31} &
\tfrac3{23} & \tfrac8{61} & \tfrac5{38} \\[2pt]
\text{count} & 1 & 17 & 4 & 87 & 275 & 4 & 17
\end{array}$\qquad
$\begin{array}[t]{lc}
\text{value} & \tfrac7{53} \;\text{and}\; \tfrac2{15} \\[2pt]
\text{count} & 21 \;\text{and}\; 1561
\end{array}$
\end{center}
The six rung values open the spectrum in exact ladder order, with the
family members first; the non-rung values that follow ($\tfrac18$,
$\tfrac7{55}$, \dots) are the next spectral strata.

\section{Machine verification}\label{sec:machine}

Every theorem in this note is accompanied by an exact machine check; all
arithmetic is integer or rational, never floating point.

\begin{enumerate}[label=\textbf{M\arabic*.}, leftmargin=*]
\item \textbf{Rung family, $n = 2..16$, $k = 1..12$} ($180$ parameter
pairs): the gap value equals $\tfrac{k}{kn+1}$; the argmax set equals
$\{\tfrac{k}{kn+1}, \tfrac{kn+1-k}{kn+1}\}$ exactly for $k \ge 2$
(including all gcd-degenerate pairs such as $n = k = 3$) and the units grid
for $k = 1$; binders are $\{1, kn\}$; every other runner sits at distance
$\ge \tfrac{k+1}{kn+1}$. Zero failures.
\item \textbf{Filler characterization}: agreement with brute force on all
$612$ zoo cores ($n = 5..8$), $480$ random cores, and the degenerate
regimes, as summarized in Section~\ref{sec:filler}.
\item \textbf{Census validation}: the four validation layers of
Section~\ref{sec:census} (identity with the unfiltered solver at two
radii, byte-level shard identity, $5{,}384/5{,}384$ independent sample
re-verification, exact accounting of all $536{,}878{,}650$ vectors).
\item \textbf{Earlier checks carried over}: the $k = 2$ case of
Theorem~\ref{thm:rung} (value, unique argmax, binders, slack) was
previously machine-verified for $n = 2..40$ with two independent proofs;
the present $k \ge 3$ extension and the removal of the gcd condition are
new here.
\end{enumerate}

\section{Open problems}\label{sec:open}

\begin{enumerate}[label=\textbf{O\arabic*.}, leftmargin=*]
\item \textbf{The $n = 8$ inter-rung theorem.} The window
$\bigl(\tfrac19, \tfrac{2}{17}\bigr)$ is empty through $V = 50$
($5.4 \times 10^8$ sets), whereas $n = 6$ and $n = 7$ have inter-rung
populations at $3/23$ and $3/23$-type values living on grids unrelated to
$2n+1$. Is the window empty at $n = 8$ for \emph{all} radii? A structural
proof --- for instance via the binding-pair lemma with modulus $17$ ---
would establish the ``second-rung conjecture'' at $n = 8$ and would be the
first such closure beyond the cases where it is known to fail.
\item \textbf{Rung-zoo structure.} At every radius tested, the rung-$k$ zoo
contains the family member $\{1, \dots, n-1, kn\}$ and decomposes into
primitive classes generated from lower-rung cores by fillers. The filler
theorem reduces this to a finite check per core; a general structural
theorem describing the zoos of the rung family remains open.
\item \textbf{$n = 9$ predictions.} The mechanism predicts: tight zoo
$\{1, \dots, 9\} \times c$; rung values $\tfrac{k}{9k+1} =
\tfrac1{10}, \tfrac2{19}, \tfrac3{28}, \dots$; rung-2 zoo from zoo-7/zoo-8
cores plus fillers binding modulo $19$ (a prime, so every pair
$\{v, 19m - v\}$ is admissible at grid level $e = 1$); inter-rung window
$\bigl(\tfrac1{10}, \tfrac2{19}\bigr)$ populated or empty --- the
first test of O1's mechanism one rung higher.
\item \textbf{Beyond the family.} The no-wrap induction exploits the
structure $\{1, \dots, n-1\} \cup \{kn\}$ through the points
$\{0, t, \dots, nt\}$. Which other speed families admit a complete argmax
classification, and is the argmax of \emph{every} rational-valued gap
attained on a rigid grid (the generalized Lemma A)?
\end{enumerate}

\begin{thebibliography}{9}

\bibitem{wills}
J.~M. Wills,
\emph{Zwei S\"atze \"uber inhomogene diophantische Approximation von
Irrationalzahlen},
Monatshefte f\"ur Mathematik \textbf{71} (1967), 261--263.

\bibitem{cusick}
T.~W. Cusick,
\emph{View-obstruction problems in $n$-dimensional cubes},
Discrete Mathematics \textbf{27} (1979), 25--28.

\bibitem{bhk}
T.~Bohman, R.~Holzman, D.~Kleitman,
\emph{Six lonely runners},
Israel Journal of Mathematics \textbf{133} (2003), 1--10.

\bibitem{repo}
MIKEAA2020,
\emph{lonely-runner: exhaustive exact checks and audit trail},
software repository and audits (2025--2026), in particular the
computational-check directory and the audit ``rung-2 ladder and R1''.

\end{thebibliography}

\end{document}
